\documentclass[aspectratio=169,11pt]{beamer}
% Shared Beamer style for practice contest solution walkthroughs.
\usepackage[utf8]{inputenc}
\usepackage[T1]{fontenc}
\usepackage{lmodern}
\usepackage{amsmath,amssymb}
\usepackage{xcolor}
\usepackage{hyperref}
\usepackage{listings}

\definecolor{CPNavy}{HTML}{172033}
\definecolor{CPBlue}{HTML}{2563EB}
\definecolor{CPGreen}{HTML}{16A34A}
\definecolor{CPOrange}{HTML}{F97316}
\definecolor{CPGray}{HTML}{64748B}
\definecolor{CPLight}{HTML}{F8FAFC}
\definecolor{CPCodeBg}{HTML}{F1F5F9}
\definecolor{CPCodeRule}{HTML}{CBD5E1}

\usetheme{default}
\usefonttheme{professionalfonts}
\setbeamertemplate{navigation symbols}{}
\setbeamersize{text margin left=0.65cm,text margin right=0.65cm}
\setbeamertemplate{itemize item}{\color{CPBlue}$\blacktriangleright$}
\setbeamertemplate{itemize subitem}{\color{CPGray}$\triangleright$}

\setbeamercolor{normal text}{fg=CPNavy,bg=white}
\setbeamercolor{frametitle}{fg=white,bg=CPNavy}
\setbeamercolor{title}{fg=white,bg=CPNavy}
\setbeamercolor{subtitle}{fg=CPLight,bg=CPNavy}
\hypersetup{colorlinks=true,urlcolor=CPBlue}

\setbeamertemplate{frametitle}{%
  \nointerlineskip%
  \begin{beamercolorbox}[wd=\paperwidth,ht=0.88cm,dp=0.22cm,leftskip=0.55cm,rightskip=0.35cm]{frametitle}
    \usebeamerfont{frametitle}\insertframetitle\hfill\small\insertframenumber/\inserttotalframenumber
  \end{beamercolorbox}%
}

\setbeamertemplate{footline}{%
  \leavevmode%
  \hbox{%
    \begin{beamercolorbox}[wd=.58\paperwidth,ht=2.4ex,dp=1ex,leftskip=.5cm]{author in head/foot}%
      \color{CPGray}\insertshorttitle
    \end{beamercolorbox}%
    \begin{beamercolorbox}[wd=.42\paperwidth,ht=2.4ex,dp=1ex,rightskip=.5cm plus1fil]{date in head/foot}%
      \hfill\color{CPGray}\insertshortauthor
    \end{beamercolorbox}%
  }%
}

\setbeamertemplate{title page}{%
  \vspace*{1.25cm}
  \begin{beamercolorbox}[wd=\paperwidth,sep=0.65cm,leftskip=0.15cm,rightskip=0.15cm]{title}
    {\color{white}\Huge\bfseries\inserttitle\par}
    \vspace{0.35cm}
    {\color{CPLight}\Large\insertsubtitle\par}
    \vspace{0.6cm}
    {\color{CPLight}\large\insertauthor\par}
  \end{beamercolorbox}
  \vfill
}

\newcommand{\sectiontag}[1]{\textbf{\color{CPBlue}#1}}
\newenvironment{tightitemize}{\begin{itemize}\small\setlength{\itemsep}{0.18cm}}{\end{itemize}}
\lstdefinestyle{cpcode}{
  language=Python,
  basicstyle=\ttfamily\tiny,
  keywordstyle=\color{CPBlue}\bfseries,
  commentstyle=\color{CPGray}\itshape,
  stringstyle=\color{CPGreen},
  backgroundcolor=\color{CPCodeBg},
  frame=single,
  framerule=0.25pt,
  rulecolor=\color{CPCodeRule},
  showstringspaces=false,
  columns=fullflexible,
  keepspaces=true,
  breaklines=true,
  breakatwhitespace=false,
  tabsize=4,
  xleftmargin=0.08cm,
  xrightmargin=0.08cm,
  aboveskip=0.08cm,
  belowskip=0.08cm
}


\title[smallestmultiple]{Smallest Multiple}
\subtitle{Day 5 solution walkthrough}
\author[Solution Deck]{Competitive Programming Bootcamp}
\date{}

\begin{document}

\begin{frame}[plain]
  \titlepage
\end{frame}

% BEGIN GENERATED STATEMENT INTERPRETATION
\begin{frame}{Interpret the Word Problem}
\small
\sectiontag{Statement excerpts:}
\begin{tightitemize}
  \item \textit{``smallest number that is a multiple of every number''}
  \item \textit{``Input consists of up to 500 test cases''}
\end{tightitemize}

\vspace{0.18cm}
\sectiontag{Programming interpretation:}
\begin{tightitemize}
  \item For each line, compute the least common multiple of all listed values.
  \item The pairwise formula uses gcd(a, b) * lcm(a, b) = a * b.
  \item Fold the list from left to right until every number has been included.
\end{tightitemize}
\end{frame}
% END GENERATED STATEMENT INTERPRETATION







\begin{frame}{Problem Model}
\small
\sectiontag{Textbook connection:} CP4 5.3.6 GCD and LCM

\vspace{0.25cm}
\sectiontag{Core technique:} Repeated LCM using GCD

\vspace{0.25cm}
\begin{tightitemize}
  \item Each input line is an independent test case.
  \item The task is to find the smallest positive integer divisible by every number on the line.
  \item Input continues until EOF.
\end{tightitemize}
\end{frame}

\begin{frame}{How to Recognize the Approach}
\begin{tightitemize}
  \item The smallest number divisible by two integers is their LCM.
  \item The LCM operation can be folded across a list of numbers.
  \item GCD gives a safe and direct way to compute each pairwise LCM.
\end{tightitemize}
\end{frame}

\begin{frame}{Algorithm}
\begin{tightitemize}
  \item Define lcm(a, b) = (a * b) // gcd(a, b).
  \item For each nonempty input line, parse all integers.
  \item Set answer to the first number.
  \item For every remaining number x, set answer = lcm(answer, x).
  \item Print the final answer.
\end{tightitemize}
\end{frame}

\begin{frame}{Walkthrough of the Key State}
\begin{tightitemize}
  \item For 2 3 4, lcm(2, 3) = 6 and lcm(6, 4) = 12.
  \item Associativity lets us combine values one at a time.
  \item Python's gcd uses the Euclidean algorithm.
\end{tightitemize}
\end{frame}

\begin{frame}{Why the Algorithm Is Correct}
\begin{tightitemize}
  \item At each step, answer is divisible by all processed numbers.
  \item The pairwise LCM is the smallest number that preserves divisibility by the previous answer and the new number.
  \item After all numbers are processed, answer is the smallest common multiple of the whole line.
\end{tightitemize}
\end{frame}

\begin{frame}{Complexity and Implementation Notes}
\small
\sectiontag{Complexity:} O(K log M) time per line for K numbers of size about M; O(K) memory to parse the line.

\vspace{0.3cm}
\begin{tightitemize}
  \item In fixed-width languages, divide before multiplying to reduce overflow risk.
  \item Empty lines are skipped.
\end{tightitemize}
\end{frame}



% BEGIN GENERATED CODE WALKTHROUGH
\begin{frame}[fragile]{Code: LCM Helper}
\scriptsize
\sectiontag{Source lines:} \texttt{smallestmultiple.py:43--51}

\vspace{0.08cm}
\begin{columns}[T,totalwidth=\textwidth]
\begin{column}{0.58\textwidth}
\begin{lstlisting}[style=cpcode]
import sys
from math import gcd

def lcm(a, b):
    return (a * b) // gcd(a, b)
\end{lstlisting}
\end{column}
\begin{column}{0.38\textwidth}
\sectiontag{What this chunk does}
\begin{tightitemize}
  \item Python's gcd gives the greatest common divisor efficiently.
  \item The helper rearranges the gcd-lcm identity to compute lcm(a, b).
  \item Integer division keeps the answer exact.
\end{tightitemize}
\end{column}
\end{columns}
\end{frame}

\begin{frame}[fragile]{Code: Read One Line at a Time}
\scriptsize
\sectiontag{Source lines:} \texttt{smallestmultiple.py:54--68}

\vspace{0.08cm}
\begin{columns}[T,totalwidth=\textwidth]
\begin{column}{0.58\textwidth}
\begin{lstlisting}[style=cpcode]
def main():
    for line in sys.stdin:

        numbers = list(map(int, line.split()))

        if not numbers:
            continue

        answer = numbers[0]
\end{lstlisting}
\end{column}
\begin{column}{0.38\textwidth}
\sectiontag{What this chunk does}
\begin{tightitemize}
  \item Each input line is one independent test case.
  \item The loop over sys.stdin naturally stops at EOF.
  \item The first number initializes the running LCM.
\end{tightitemize}
\end{column}
\end{columns}
\end{frame}

\begin{frame}[fragile]{Code: Fold the LCM}
\scriptsize
\sectiontag{Source lines:} \texttt{smallestmultiple.py:69--82}

\vspace{0.08cm}
\begin{columns}[T,totalwidth=\textwidth]
\begin{column}{0.58\textwidth}
\begin{lstlisting}[style=cpcode]
        for i in range(1, len(numbers)):
            answer = lcm(answer, numbers[i])

        print(answer)

main()
\end{lstlisting}
\end{column}
\begin{column}{0.38\textwidth}
\sectiontag{What this chunk does}
\begin{tightitemize}
  \item Each remaining number is combined with the current answer.
  \item After the loop, answer is divisible by every number on the line.
  \item The result is printed immediately for that test case.
\end{tightitemize}
\end{column}
\end{columns}
\end{frame}
% END GENERATED CODE WALKTHROUGH

\begin{frame}{Source}
\small
\sectiontag{Solution file:} \texttt{practice\_contests/day\_5/smallestmultiple.py}

\vspace{0.3cm}
\sectiontag{Problem:} \url{https://open.kattis.com/problems/smallestmultiple}

\vspace{0.45cm}
Use this deck with the source code open beside it. The slides explain the reasoning; the code shows the exact input handling and output formatting.
\end{frame}

\end{document}
